\documentclass[conference]{IEEEtran}
\IEEEoverridecommandlockouts

\usepackage[utf8]{inputenc}
\usepackage[T1]{fontenc}
\usepackage{cite}
\usepackage{amsmath,amssymb,amsfonts}
\usepackage{algorithmic}
\usepackage{graphicx}
\usepackage{textcomp}
\usepackage{xcolor}
\usepackage{booktabs}
\usepackage{multirow}
\usepackage{url}
\usepackage{hyperref}

% ----------------------------------------------------------------
% TODO: substituir [AUTHOR] e [INSTITUTION] antes de submeter
% ----------------------------------------------------------------

\begin{document}

\title{A Generic TinyML Pipeline for Time-Series Anomaly Detection:\\
Edge-Aware Preprocessing, MLOps Quality Gates,\\
and Secure OTA Updates on ESP32}

\author{
  \IEEEauthorblockN{[AUTHOR NAME]}
  \IEEEauthorblockA{
    \textit{[Department]}\\
    \textit{[Institution]}\\
    [City, Country]\\
    [email]
  }
}

\maketitle

% ----------------------------------------------------------------
\begin{abstract}
We present a generic, reproducible TinyML pipeline for anomaly detection
in complex time series, validated in the seismic domain targeting
ESP32 microcontrollers.
The architecture decouples domain adapters from the machine learning core,
enabling reuse for industrial vibration, electrical current, or telemetry
sensors without pipeline rewriting.
Edge-aware preprocessing eliminates MCU-incompatible steps
(instrumental response removal), reducing training-serving skew at
embedded inference time.
The pipeline compares six model families---from STA/LTA to
Optuna-optimized dilated convolutional networks---under a
comprehensive set of operational metrics: AUC-PR, VUS-PR,
PA-F1, event-level F1, and false positives per hour (FP/h).
The best model, a Tiny TCN with 14{,}897 parameters, achieves AUC-PR of
0.9416 and 3.136 FP/h, outperforming all baselines including the
traditional STA/LTA method.
A complete MLOps cycle---DVC, MLflow, multidimensional quality gate,
model promotion, and an OTA package with SHA-256 integrity verification
and rollback---establishes an auditable foundation for production TinyML systems.
The full pipeline and all trained models are publicly available.
\end{abstract}

\begin{IEEEkeywords}
TinyML, anomaly detection, time series, seismology, ESP32, MLOps,
temporal convolutional network, over-the-air update.
\end{IEEEkeywords}

% ================================================================
\section{Introduction}
\label{sec:intro}
% ================================================================

Continuous sensor monitoring systems face a recurring operational dilemma:
transmitting all sensor data for centralized analysis incurs high energy,
bandwidth, and storage costs, while fixed-threshold alerting on the device
produces high false-positive rates and poor adaptability~\cite{sculley2015hidden}.
In domains such as seismic monitoring, industrial vibration, and equipment
telemetry, continuous acquisition generates data volumes that make centralized
analysis impractical at scale.

TinyML---the execution of machine learning models on microcontrollers with
severely limited resources---offers an alternative: the edge device processes
local signal windows and forwards data only when there is reliable evidence
of an anomaly~\cite{warden2019tinyml,banbury2021mlperf}.
This reduces transmission, energy consumption, and operational cost, while
enabling faster response and offline operation.

In the seismic domain, automatic event detection has been addressed by classical
methods such as STA/LTA~\cite{allen1978automatic,withers1998comparison} and,
more recently, deep networks such as PhaseNet~\cite{zhu2019phasenet} and
EQTransformer~\cite{mousavi2020earthquake}.
These models, however, are designed for server execution and do not consider
the memory and processing constraints of microcontrollers.

Evaluation methodology is also under scrutiny:
Kim et al.~\cite{kim2022towards} show that standard AUC-PR is unreliable
for windowed time series due to point-to-point alignment sensitivity,
proposing VUS-PR as a more robust alternative.
Xu et al.~\cite{xu2018donut} and Tatbul et al.~\cite{tatbul2018precision}
formalize point-adjust F1 (PA-F1) and event-level F1 as operationally
meaningful complements.

This paper contributes:
\begin{enumerate}
  \item A \emph{generic, reproducible} TinyML pipeline for binary time-series
        anomaly detection, domain-agnostic via adapter separation.
  \item \emph{Edge-aware preprocessing} that eliminates MCU-incompatible
        steps, reducing training-serving skew.
  \item Comparison of six model families under a \emph{full operational
        metric set}: AUC-PR, VUS-PR, PA-F1, Event-F1, and FP/h.
  \item A \emph{complete MLOps cycle}: DVC, MLflow, multidimensional quality
        gate, model promotion, and an OTA package with SHA-256 integrity
        verification and simulated rollback.
\end{enumerate}

% ================================================================
\section{Related Work}
\label{sec:related}
% ================================================================

\subsection{TinyML and Embedded Inference}

Warden and Situnayake~\cite{warden2019tinyml} established practical foundations
for running ML models on microcontrollers.
Banbury et al.~\cite{banbury2021mlperf} formalized performance benchmarks for
TinyML systems, exposing tradeoffs between accuracy, latency, and energy.
David et al.~\cite{david2021tensorflow} present TensorFlow Lite Micro as
the reference framework for OS-free inference.
Lin et al.~\cite{lin2020mcunet} propose architecture and inference engine
co-design to maximize accuracy on MCUs with less than 512~KB SRAM.

\subsection{Anomaly Detection in Time Series}

Chandola et al.~\cite{chandola2009anomaly} provide a comprehensive taxonomy
of anomaly detection techniques.
Pang et al.~\cite{pang2021deep} review deep learning methods for anomalies,
highlighting the importance of appropriate metrics for imbalanced distributions.
Kim et al.~\cite{kim2022towards} propose VUS-PR, integrating AUC-PR over
multiple prediction buffer sizes to produce a more rigorous evaluation measure.
Tatbul et al.~\cite{tatbul2018precision} define range-based precision and
recall for time series, formalizing event-level evaluation.

\subsection{Seismic Detection with ML}

STA/LTA~\cite{allen1978automatic,withers1998comparison} is the classical
baseline: it compares short-window energy (STA) to long-window energy (LTA)
to identify event onsets.
PhaseNet~\cite{zhu2019phasenet} is a U-Net for seismic phase detection at
datacenter scale.
EQTransformer~\cite{mousavi2020earthquake} introduces a hierarchical
transformer for simultaneous detection and phase classification.
Both achieve strong performance but with tens of millions of parameters,
infeasible for ESP32 deployment.

\subsection{Temporal Convolutional Networks}

Bai et al.~\cite{bai2018empirical} demonstrate empirically that causal dilated
convolutional networks (TCN) outperform LSTMs on diverse sequence modeling tasks.
The absence of recurrent state simplifies sequential inference and reduces
memory consumption, an advantage relevant for embedded devices.

\subsection{MLOps for Embedded Systems}

Sculley et al.~\cite{sculley2015hidden} identify hidden technical debt in ML
systems, including training-serving skew and lack of data versioning.
Amershi et al.~\cite{amershi2019software} document software engineering
practices for production ML.
DVC~\cite{dvc} and MLflow~\cite{mlflow} form the recognized basis for
reproducible MLOps pipelines.

% ================================================================
\section{Methodology}
\label{sec:method}
% ================================================================

\subsection{Pipeline Architecture}

The pipeline enforces explicit separation between domain adapters and the
ML core (Fig.~\ref{fig:pipeline}):

\begin{figure}[htbp]
  \centering
  % TODO: inserir figura do pipeline
  % \includegraphics[width=\columnwidth]{figures/pipeline.pdf}
  \framebox[\columnwidth]{\parbox{0.9\columnwidth}{\centering
    \texttt{raw data $\to$ domain adapter $\to$ generic NPZ dataset} \\
    \texttt{$\to$ contract validation $\to$ feature extraction / raw signal} \\
    \texttt{$\to$ model training $\to$ Optuna HPO $\to$ candidate manifest} \\
    \texttt{$\to$ quality gate $\to$ production manifest} \\
    \texttt{$\to$ TFLite export + C header $\to$ OTA package}
  }}
  \caption{End-to-end TinyML pipeline. Domain adapters are the only
           domain-specific components; everything downstream is generic.}
  \label{fig:pipeline}
\end{figure}

The \emph{dataset contract} is a NumPy compressed archive with keys
\texttt{X\_\{train,val,test\}} of shape \texttt{(N, window\_size)} for
univariate series or \texttt{(N, window\_size, n\_channels)} for
multivariate series, and binary labels \texttt{y\_\{train,val,test\}}.
Adding a new domain requires only a new adapter that respects this contract.

\subsection{Seismic Dataset}
\label{sec:dataset}

% TODO: preencher com descrição real dos dados
The validation domain is seismic.
Raw data are MiniSEED files organized as:
\texttt{raw/events/} (anomalous windows, $y=1$) and
\texttt{raw/continuous/} (background noise, $y=0$).

\begin{table}[htbp]
  \caption{Seismic Profile Parameters (\texttt{seismic\_edge\_v1})}
  \label{tab:profile}
  \centering
  \begin{tabular}{ll}
    \toprule
    Parameter & Value \\
    \midrule
    Sampling rate     & 40~Hz \\
    Window size       & 800 samples (20~s) \\
    Window step       & 10~s (50\% overlap) \\
    Split strategy    & By event \\
    Test anomaly ratio & $\approx$12.8\% \\
    \bottomrule
  \end{tabular}
\end{table}

Split by event guarantees that windows from the same earthquake do not
appear simultaneously in train and test, preventing data leakage.

\subsection{Edge-Aware Preprocessing}

\textbf{Offline (training):}
\texttt{resample 40~Hz $\to$ linear detrend $\to$ demean $\to$ 5\% taper
$\to$ bandpass 0.5--15~Hz $\to$ per-window z-score}.

\textbf{Edge (ESP32 inference):}
\texttt{linear detrend $\to$ 5\% taper $\to$ bandpass 0.5--15~Hz
$\to$ per-window z-score}.

The \texttt{remove\_response} step (requiring StationXML) was removed---it
cannot be reproduced on the MCU.
All retained steps are implementable in C/C++ with single-precision
floating-point arithmetic.

\subsection{Feature Extraction}

For classical supervised models, 28 features are extracted per window:

\textbf{Temporal (20):} mean, std, min, max, median, abs\_mean, abs\_peak,
RMS, crest factor, peak-to-peak, energy, skewness, kurtosis, $p_{05}$,
$p_{25}$, $p_{75}$, $p_{95}$, IQR, zero crossings, zero-crossing rate.

\textbf{Spectral (8):} dominant frequency, spectral centroid,
spectral rolloff (85\%), bandpower in 0--3, 0.5--3, 3--8, 8--15~Hz,
spectral entropy.

Neural networks (Tiny CNN, Tiny TCN) operate directly on 800 raw samples.

\subsection{Models}

\begin{table}[htbp]
  \caption{Model Families Evaluated}
  \label{tab:models}
  \centering
  \begin{tabular}{llcc}
    \toprule
    Family & Model & Input & Edge \\
    \midrule
    Classical baseline & STA/LTA              & Raw signal   & \checkmark \\
    Linear             & Logistic Regression  & 28 features  & $\times$   \\
    Ensemble           & Random Forest        & 28 features  & $\times$   \\
    Ensemble           & Extra Trees          & 28 features  & $\times$   \\
    Lightweight neural & Tiny CNN             & 800 pts      & \checkmark \\
    Lightweight neural & Tiny TCN (Optuna)    & 800 pts      & \checkmark \\
    \bottomrule
  \end{tabular}
\end{table}

\textbf{Tiny TCN (best model):}
Causal dilated separable Conv1D blocks with residual connections.
Final Optuna configuration (trial~24, 88 trials total):
6 blocks, 32 filters, kernel 7, dilation base 3, dense 24,
dropout 0.028, spatial dropout 0.066; \textbf{14{,}897 parameters}.

\subsection{Hyperparameter Optimization}

Optuna~\cite{akiba2019optuna} maximizes a composite objective:

\begin{equation}
  \mathcal{J} = \text{AUC-PR}_\text{val} - \lambda \cdot \max(0,\; \text{FP/h} - \text{FP/h}_\text{tol})
  \label{eq:objective}
\end{equation}

\noindent where $\lambda = 0.01$ and $\text{FP/h}_\text{tol} = 10$.
The decision threshold is chosen on the validation set and applied
unchanged on the test set.

\subsection{Evaluation Metrics}

\textbf{AUC-PR:} primary metric for imbalanced datasets.

\textbf{VUS-PR}~\cite{kim2022towards}:
average AUC-PR over prediction buffers $b \in [0, b_{\max}]$:
\begin{equation}
  \text{VUS-PR} = \frac{1}{b_{\max}+1} \sum_{b=0}^{b_{\max}}
  \text{AUC-PR}^{(b)}
  \label{eq:vuspr}
\end{equation}
where $\text{AUC-PR}^{(b)}$ dilates each anomaly segment by $b$ windows.
Buffer $b=0$ recovers standard AUC-PR.

\textbf{PA-F1}~\cite{xu2018donut}:
point-adjust F1---if any window within an anomaly segment is correctly
detected, all windows in that segment are counted as TP.

\textbf{Event-F1}~\cite{tatbul2018precision}:
F1 at segment level---a true segment is a TP if $\geq10\%$ of its
windows are predicted anomalous.

\textbf{FP/h:} operational false positive rate per monitoring hour:
\begin{equation}
  \text{FP/h} = \frac{\text{FP}}{\text{hours of evaluated signal}}
  \label{eq:fph}
\end{equation}

\subsection{MLOps Cycle}

\textbf{DVC} versions data and pipeline stages: \texttt{generate\_data $\to$
validate\_dataset $\to$ train\_all $\to$ export\_tflite}.

\textbf{MLflow} tracks all experiment parameters, metrics, and model artifacts.

\textbf{Quality gate} enforces multidimensional thresholds before promotion:
AUC-PR $\geq 0.80$, FP/h $\leq 10$, $|\Delta\text{AUC-PR}_\text{val/test}| \leq 0.05$,
model size $\leq 100$~KB.

\textbf{OTA package} contains: versioned TFLite model, JSON manifest,
SHA-256 digest, compatibility metadata (target, runtime, preprocessing version),
and a simulated rollback mechanism.

% ================================================================
\section{Results}
\label{sec:results}
% ================================================================

\subsection{Model Comparison}

\begin{table*}[htbp]
  \caption{Model Comparison on the Seismic Test Set}
  \label{tab:results}
  \centering
  \begin{tabular}{lcccccccc}
    \toprule
    Model & AUC-PR & VUS-PR & PA-F1 & Event-F1 & F1 & FP/h & Params & Edge \\
    \midrule
    \textbf{Optuna Tiny TCN}
          & \textbf{0.9416} & \textbf{[TODO]} & \textbf{[TODO]}
          & \textbf{[TODO]} & \textbf{0.8885} & \textbf{3.136}
          & \textbf{14{,}897} & \checkmark \\
    Optuna Tiny CNN
          & 0.9127 & [TODO] & [TODO] & [TODO] & 0.8526 & 4.896
          & 19{,}489 & \checkmark \\
    Tiny TCN (no HPO)
          & 0.8964 & [TODO] & [TODO] & [TODO] & 0.7666 & 21.984
          & $\approx$6k & \checkmark \\
    Optuna Random Forest
          & 0.8127 & [TODO] & [TODO] & [TODO] & 0.7367 & 9.264
          & --- & $\times$ \\
    Optuna Extra Trees
          & 0.7901 & [TODO] & [TODO] & [TODO] & 0.7102 & 11.296
          & --- & $\times$ \\
    STA/LTA
          & 0.1662 & [TODO] & [TODO] & [TODO] & 0.2760 & ---
          & --- & \checkmark \\
    \bottomrule
  \end{tabular}
  % TODO: preencher VUS-PR, PA-F1, Event-F1 rodando evaluate_scores
  %       nos scores salvos no MLflow (sem retreinamento)
\end{table*}

The Optuna Tiny TCN dominates across all metrics: highest AUC-PR (0.9416),
lowest FP/h (3.136) among neural models, and only 14{,}897 parameters---well
within ESP32 flash capacity after TFLite conversion.

STA/LTA achieves AUC-PR of 0.1662, near random baseline for the class
distribution, confirming that fixed-threshold methods are insufficient
for the variability of the evaluated data.

\subsection{Impact of Hyperparameter Optimization}

\begin{table}[htbp]
  \caption{HPO Convergence for Tiny TCN}
  \label{tab:hpo}
  \centering
  \begin{tabular}{lcc}
    \toprule
    Configuration & AUC-PR (val) & FP/h \\
    \midrule
    No HPO          & ---    & 21.984 \\
    Trial 0         & 0.8152 & 8.539  \\
    Trial 8         & 0.8906 & 4.521  \\
    Trial 11        & 0.8981 & 3.586  \\
    \textbf{Trial 24 (best)} & \textbf{0.9343} & \textbf{3.136} \\
    \bottomrule
  \end{tabular}
\end{table}

HPO reduces FP/h by \textbf{85.7\%} relative to the default architecture
(21.984 $\to$ 3.136).
The improvement in AUC-PR from validation to test ($+0.0079$) indicates
good generalization with no overfitting.

\subsection{Feature Importance (Classical Models)}

Feature importance analysis on Random Forest reveals that spectral features
dominate class separation:
\texttt{spectral\_rolloff\_85}, \texttt{bandpower\_8\_15hz},
\texttt{spectral\_centroid}, \texttt{zero\_crossing\_rate},
\texttt{kurtosis}, \texttt{spectral\_entropy}.
Anomalies are discriminated by spectral distribution and waveform shape,
not by raw amplitude---consistent with seismic physics.

\subsection{MLOps Validation}

The Tiny TCN passed the quality gate, generating \texttt{production\_manifest.json}
and an OTA package (\texttt{seismic\_edge\_v1\_tiny\_tcn\_v4-2}) with verified
SHA-256 digest.
The simulated OTA flow correctly detected available updates, applied the new
version, suppressed unnecessary reinstallation, and completed rollback
after a simulated failure---all state transitions verified.

\subsection{Embedded Resource Estimates}

\begin{table}[htbp]
  \caption{Embedded Resource Estimates for ESP32}
  \label{tab:resources}
  \centering
  \begin{tabular}{lc}
    \toprule
    Artifact & Size \\
    \midrule
    \texttt{model.tflite} (float32) & [TODO] KB \\
    \texttt{model.tflite} (int8)    & [TODO] KB \\
    \texttt{model\_data.h}          & [TODO] KB \\
    Estimated inference RAM         & [TODO] KB \\
    Estimated inference latency     & [TODO] ms \\
    \bottomrule
  \end{tabular}
  % TODO: medir após exportação TFLite com quantização int8
\end{table}

% ================================================================
\section{Discussion}
\label{sec:discussion}
% ================================================================

\textbf{Edge-aware preprocessing.}
Removing \texttt{remove\_response} from the training pipeline aligns the
offline and embedded execution paths, reducing the most common source of
training-serving skew in physics-based signal processing.
Residual risk: numerical differences between the Python and C implementations
of bandpass filtering.
Hardware validation on ESP32 (Section~\ref{sec:conclusion}, future work)
will quantify this residual gap.

\textbf{VUS-PR vs.\ AUC-PR.}
With 50\% window overlap, the same anomaly contributes to multiple positive
windows.
AUC-PR may overestimate performance by counting a single missed event as
multiple false negatives, inflating recall.
VUS-PR provides a more conservative and operationally meaningful estimate;
values will be populated upon running the evaluation pipeline on stored scores.

\textbf{Compact models for edge deployment.}
The Tiny TCN (14{,}897 params) and Tiny CNN (19{,}489 params) are both
viable for ESP32 flash.
Non-edge models (Random Forest, Extra Trees) require megabytes after
joblib serialization, precluding MCU deployment.

% ================================================================
\section{Conclusion}
\label{sec:conclusion}
% ================================================================

We presented a generic, reproducible TinyML pipeline covering the complete
cycle from raw data acquisition to secure OTA model updates on ESP32.
The architecture decouples domain adapters from the ML core, enabling
reuse across sensor domains without pipeline rewriting.
Edge-aware preprocessing eliminates MCU-incompatible dependencies, reducing
training-serving skew.
The Optuna-optimized Tiny TCN achieves AUC-PR of 0.9416 and 3.136 FP/h
with 14{,}897 parameters, outperforming all evaluated baselines.
The multidimensional quality gate and SHA-256-verified OTA package with
rollback establish an auditable foundation for production TinyML systems.

Future work includes:
(i) hardware validation on ESP32 to quantify the preprocessing gap;
(ii) cross-domain replication on Petrobras 3W and C-MAPSS datasets
     (see companion papers);
(iii) int8 quantization and on-device latency measurement.

% ================================================================
% REFERENCES
% ================================================================

\begin{thebibliography}{00}

\bibitem{sculley2015hidden}
D.~Sculley et al.,
``Hidden technical debt in machine learning systems,''
in \textit{Proc. NeurIPS}, 2015, pp. 2503--2511.

\bibitem{warden2019tinyml}
P.~Warden and D.~Situnayake,
\textit{TinyML: Machine Learning with TensorFlow Lite on Arduino and
Ultra-Low-Power Microcontrollers}.
O'Reilly Media, 2019.

\bibitem{banbury2021mlperf}
C.~Banbury et al.,
``MLPerf Tiny Benchmark,''
in \textit{Proc. MLSys}, 2021.

\bibitem{david2021tensorflow}
R.~David et al.,
``TensorFlow Lite Micro: Embedded machine learning for TinyML systems,''
in \textit{Proc. MLSys}, 2021.

\bibitem{lin2020mcunet}
J.~Lin et al.,
``MCUNet: Tiny deep learning on IoT devices,''
in \textit{Proc. NeurIPS}, 2020.

\bibitem{allen1978automatic}
R.~Allen,
``Automatic earthquake recognition and timing from single traces,''
\textit{Bull. Seismol. Soc. Am.}, vol.~68, no.~5, pp. 1521--1532, 1978.

\bibitem{withers1998comparison}
M.~Withers et al.,
``A comparison of select trigger algorithms for automated global seismic
phase and event detection,''
\textit{Bull. Seismol. Soc. Am.}, vol.~88, no.~1, pp. 95--106, 1998.

\bibitem{zhu2019phasenet}
W.~Zhu and G.~C.~Beroza,
``PhaseNet: A deep-neural-network-based seismic arrival-time picking method,''
\textit{Geophys. J. Int.}, vol.~216, no.~1, pp. 261--273, 2019.

\bibitem{mousavi2020earthquake}
S.~M.~Mousavi et al.,
``Earthquake transformer---an attentive deep-learning model for simultaneous
earthquake detection and phase picking,''
\textit{Nat. Commun.}, vol.~11, p. 3952, 2020.

\bibitem{bai2018empirical}
S.~Bai, J.~Z.~Kolter, and V.~Koltun,
``An empirical evaluation of generic convolutional and recurrent networks
for sequence modeling,''
arXiv:1803.01271, 2018.

\bibitem{chandola2009anomaly}
V.~Chandola, A.~Banerjee, and V.~Kumar,
``Anomaly detection: A survey,''
\textit{ACM Comput. Surv.}, vol.~41, no.~3, pp. 1--58, 2009.

\bibitem{pang2021deep}
G.~Pang et al.,
``Deep learning for anomaly detection: A review,''
\textit{ACM Comput. Surv.}, vol.~54, no.~2, pp. 1--38, 2021.

\bibitem{kim2022towards}
S.~Kim et al.,
``Towards a rigorous evaluation of time-series anomaly detection,''
in \textit{Proc. AAAI}, 2022.

\bibitem{xu2018donut}
H.~Xu et al.,
``Unsupervised anomaly detection via variational auto-encoder for seasonal
KPIs in web applications,''
in \textit{Proc. WWW}, 2018.

\bibitem{tatbul2018precision}
N.~Tatbul et al.,
``Precision and recall for time series,''
in \textit{Proc. NeurIPS}, 2018.

\bibitem{amershi2019software}
S.~Amershi et al.,
``Software engineering for machine learning: A case study,''
in \textit{Proc. ICSE-SEIP}, 2019.

\bibitem{akiba2019optuna}
T.~Akiba et al.,
``Optuna: A next-generation hyperparameter optimization framework,''
in \textit{Proc. KDD}, 2019.

\bibitem{dvc}
{Iterative AI},
``DVC: Data version control,''
\url{https://dvc.org}, 2024.

\bibitem{mlflow}
M.~Zaharia et al.,
``Accelerating the machine learning lifecycle with MLflow,''
\textit{IEEE Data Eng. Bull.}, vol.~41, no.~4, 2018.

\end{thebibliography}

\end{document}
